\subsubsection{Modulus refinement and preservation of the two readers}
\label{mg03:refinamiento}

Identity \eqref{eq:defecto-integrado} preserves the integer defect
\(1215\) through residue \(54\) and quotient \(129\). Its
generalization requires separate specifications of the reading modulus,
the support size and the number of coordinates. These three operations
produce distinct counting laws from the same APP sheets.

For \(N\geq2\), fix the positive representative
\[
 \rho_N(s)=1+((s-1)\bmod N)\in\{1,\ldots,N\},
 \qquad s=\rho_N(s)+NQ_N(s),\quad Q_N(s)\in\mathbb Z.
\]
On \(\{1,\ldots,N\}^2\), the full-modulus readings are
\(\rho_N(i+j)\) and \(\rho_N(ij)\). Their sums are denoted by
\(S_+^{\mathrm{full}}(N)\) and \(S_\times^{\mathrm{full}}(N)\).

\begin{proposition}[Defect law for an integer modulus]
For every \(N\geq2\),
\begin{align}
 S_+^{\mathrm{full}}(N)&=\frac{N^2(N+1)}2,
 \label{mg03:suma-plena}\\
 S_\times^{\mathrm{full}}(N)&=
 \frac N2\left(N^2+\sum_{i=1}^{N}\gcd(i,N)\right),
 \label{mg03:producto-pleno}\\
 \Delta_N^{\mathrm{full}}&=
 \frac N2\left(\sum_{i=1}^{N}\gcd(i,N)-N\right).
 \label{mg03:defecto-pleno}
\end{align}
\end{proposition}

\begin{proof}
In an additive row, translation by \(i\) permutes the \(N\) residue
classes and their positive representatives. Each row sums to
\(N(N+1)/2\), proving \eqref{mg03:suma-plena}. For a multiplicative
row, let \(g=\gcd(i,N)\). The image of \(j\mapsto ij\pmod N\)
contains the \(N/g\) residue classes divisible by \(g\), and each
has \(g\) preimages. Their positive representatives are
\(g,2g,\ldots,N\), whose sum is
\[
 g\frac{(N/g)(N/g+1)}2=\frac{N(N+g)}{2g}.
\]
Multiplicity \(g\) gives \(N(N+g)/2\) per row. Summation over
\(i\) yields \eqref{mg03:producto-pleno}; subtracting
\eqref{mg03:suma-plena} gives \eqref{mg03:defecto-pleno}.
\end{proof}

\begin{corollary}[Prime-power moduli]
If \(N=p^m\), where \(p\) is prime and \(m\geq1\), then
\begin{equation}
 \sum_{i=1}^{p^m}\gcd(i,p^m)
 =p^m+m(p-1)p^{m-1},\qquad
 \Delta_{p^m}^{\mathrm{full}}
 =\frac{m(p-1)}2p^{2m-1}.
 \label{mg03:primo-potencia}
\end{equation}
\end{corollary}

\begin{proof}
For each \(0\leq k<m\), there are
\(p^{m-k}-p^{m-k-1}\) indices with exact \(p\)-adic valuation
\(k\). Their greatest common divisor with \(p^m\) is \(p^k\),
so this stratum contributes \((p-1)p^{m-1}\). The \(m\) strata
contribute \(m(p-1)p^{m-1}\), whereas the index \(i=p^m\)
contributes \(p^m\). Substitution into \eqref{mg03:defecto-pleno}
gives the second formula.
\end{proof}

The triadic specialization is
\(\Delta_{3^m}^{\mathrm{full}}=m3^{2m-1}\). The first moduli
display both sums together with their difference:
\[
\begin{array}{c|r|r|r}
 N&S_+^{\mathrm{full}}&S_\times^{\mathrm{full}}&\Delta_N^{\mathrm{full}}\\\hline
 3&18&21&3\\
 9&405&459&54\\
 27&10206&10935&729\\
 81&269001&277749&8748
\end{array}
\]

Periodic refinement instead retains \(\rho_9\) and enlarges the
support to \(\{1,\ldots,9r\}^2\), with integer \(r\geq1\). Each pair of residues modulo
nine has \(r^2\) representatives in this support, for both addition
and multiplication. Consequently,
\begin{equation}
 S_+^{\mathrm{per}}(9r;9)=r^2\,405,\quad
 S_\times^{\mathrm{per}}(9r;9)=r^2\,459,\quad
 \Delta^{\mathrm{per}}(9r;9)=r^2\,54.
 \label{mg03:periodico}
\end{equation}
For \(r=3\), the sums are \(3645\) and \(4131\), and the defect
is \(486\). The full-modulus reading at \(27\) produces \(729\)
because it updates both the modulus and its representatives; the periodic
reading produces \(486\) because it retains modulus nine and multiplies
its incidences by nine. In both readings, the divisions
\(i+j=R^++NQ^+\), \(ij=R^\times+NQ^\times\), with the modulus
actually employed, preserve the integer identity at each cell and after
summation. This specifies the datum accompanying each defect.

\subsubsection{Dimensional variation at fixed modulus nine}
\label{mg03:dimension}

A second extension retains the modulus and increases the number of
factors. For \(d\geq1\), let
\[
 \mathcal A_d=\{1,\ldots,9\}^d,\qquad
 \Sigma_d(x)=\rho_9\!\left(\sum_{j=1}^d x_j\right),\qquad
 \Pi_d(x)=\rho_9\!\left(\prod_{j=1}^d x_j\right).
\]

\begin{proposition}[Dimensional count for accumulation and multiplication]
The sums of the two readers over \(\mathcal A_d\) are
\begin{equation}
 S_{+,d}=45\,9^{d-1},\qquad
 S_{\times,d}=9^{d+1}-9(d+3)6^{d-1}.
 \label{mg03:ley-dimensional}
\end{equation}
Therefore,
\begin{equation}
 \frac{S_{\times,d}-S_{+,d}}{9^d}
 =4-(d+3)\left(\frac23\right)^{d-1}
 \longrightarrow4.
 \label{mg03:limite-dimensional}
\end{equation}
\end{proposition}

\begin{proof}
For each fixed choice of the first \(d-1\) coordinates, the final
coordinate traverses every additive residue once. Each positive
representative occurs \(9^{d-1}\) times, and their sum is \(45\);
this proves the additive formula.

For multiplication, distinguish the units \(U=\{1,2,4,5,7,8\}\),
the multiples of three of valuation one \(T=\{3,6\}\), and the
zero representative \(9\). The \(6^d\) words whose factors all
belong to \(U\) distribute their products uniformly over the six
units: fixing \(d-1\) factors leaves a permutation of the group
under multiplication by the last factor. Their contribution is
\(27\,6^{d-1}\).

The words with exactly one factor in \(T\) and all other factors
in \(U\) number \(2d6^{d-1}\). For each position of the singular
factor and each choice of the remaining units, its two possible values
\(3,6\) produce each of the residues \(3,6\) once. Their contribution
is \(9d6^{d-1}\). Every remaining word contains a factor \(9\),
or at least two factors in \(T\), so its product has representative
\(9\). Their number is \(9^d-6^d-2d6^{d-1}\). Thus
\[
 S_{\times,d}=27\,6^{d-1}+9d6^{d-1}
 +9\bigl(9^d-6^d-2d6^{d-1}\bigr),
\]
which simplifies to \eqref{mg03:ley-dimensional}. Dividing by
\(9^d\) and subtracting the additive mean \(5\) gives
\eqref{mg03:limite-dimensional}; the term
\((d+3)(2/3)^{d-1}\) tends to zero.
\end{proof}

Dimension two recovers \(S_{+,2}=405\), \(S_{\times,2}=459\)
and total defect \(54\). The limit \(4\) belongs to the mean
defect under increasing dimension. Total defect, preserved quotient
and dimensional mean are specified readers of APP arithmetic; TRIT
and TPK subsequently transport these data with their orientation,
modulus and mathematical residence.
